\documentclass[11pt,a4paper]{article}
\usepackage[margin=2.2cm]{geometry}
\usepackage[T1]{fontenc}
\usepackage{booktabs,tabularx,longtable,xcolor,enumitem}
\usepackage{hyperref}
\hypersetup{colorlinks,linkcolor=blue!50!black,urlcolor=blue!50!black}
\setlist{itemsep=1pt,topsep=2pt}
\newcommand{\code}[1]{\texttt{\small #1}}
\newcolumntype{L}{>{\raggedright\arraybackslash}X}
% One block per data kind: fields, time stamps, users, and the effect of a gap.
\newcommand{\kind}[5]{%
\subsection{#1}
\noindent\begin{tabularx}{\linewidth}{@{}p{2.6cm}L@{}}
\toprule
Fields & #2\\\midrule
Time stamps & #3\\\midrule
Used by & #4\\\midrule
When missing & #5\\\bottomrule
\end{tabularx}}

\title{Data contract of the HY Analyst}
\author{FundamentalQuantmental}
\date{23 October 2026}

\begin{document}
\sloppy
\maketitle

\noindent\fcolorbox{gray}{gray!6}{\parbox{\dimexpr\linewidth-2\fboxsep-2\fboxrule}{%
\textbf{What this is.} The checklist of every kind of data that the HY Analyst reads. Use it when you connect a new
data source on another machine.

\smallskip
\textbf{One connector.} The repo holds the method. The data plugs in through one class that implements the
protocols in \code{hy/data/source.py}.

\smallskip
\textbf{One rule.} Every input reaches the analyst through \code{DataView(source, decision\_time)}. The view returns
only what was public at the decision time.}}

\section{Rules for every kind}
\begin{enumerate}
\item Every input reaches the analyst through one \code{DataView} (\code{hy/data/view.py}).
\item Every record carries \code{as\_of} and \code{available\_at}.
  \begin{itemize}
  \item \code{as\_of} is the date or period that the value describes.
  \item \code{available\_at} is the first moment the market could read the value.
  \item \code{available\_at} is never before \code{as\_of}. The record class rejects such a record.
  \end{itemize}
\item The view filters by \code{available\_at}. It also checks what the source returns. A record public after the
  decision time raises \code{LookAheadError}.
\item A source that cannot supply an item raises \code{DataUnavailable(kind, key, reason)} (\code{hy/data/pit.py}).
  The analyst turns most of these into an open question in the memo. Section~\ref{sec:kinds} states what each gap does.
\item Units:
  \begin{itemize}
  \item Rates, yields, spreads and shares are decimals. 0.065 means 6.5\%. 0.0350 means 350 bp.
  \item Amounts are in USD millions.
  \item Tenors, terms and maturities are in years.
  \item Times are naive UTC \code{datetime} values. A filing's time is its EDGAR acceptance time.
  \end{itemize}
\item Bond observations and the panels are pandas frames, because of their size. Each frame has an
  \code{available\_at} column. All other kinds are frozen dataclasses in \code{hy/data/records.py}.
\end{enumerate}

\section{The kinds at a glance}
\noindent\begin{tabularx}{\linewidth}{@{}lLllc@{}}
\toprule
No. & Kind & Source method & View method & Required\\\midrule
3.1 & Issuer identity & \code{issuers} & \code{issuer} & yes\\
3.2 & Quarterly financials & \code{financials} & \code{financials} & yes\\
3.3 & Filing facts & \code{facts} & \code{facts} & yes\\
3.4 & Filing list & \code{filings} & \code{filings} & no\\
3.5 & Balance sheet and schedules & \code{balance\_sheet} & \code{balance\_sheet} & yes\\
3.6 & Filing text & \code{filing\_text} & \code{filing\_text} & no\\
3.7 & First-reported cash flows & \code{reported\_cash\_flow} & \code{reported\_cash\_flow} & no\\
3.8 & Equity market value & \code{market\_value} & \code{market\_value} & no\\
3.9 & Bond terms & \code{bonds} & \code{bonds} & yes\\
3.10 & Bond observations & \code{observations} & \code{observations} & yes\\
3.11 & Treasury curve & \code{curve}, \code{curves} & \code{curve}, \code{curve\_history} & yes\\
3.12 & HY index level & \code{hy\_index} & \code{hy\_index} & yes\\
3.13 & Macro forecasts & \code{forecasts}, \code{forecast\_status} & same & no\\
3.14 & Peer panels & \code{panel} & \code{panel} & no\\
3.15 & Classifications, analyst inputs & \code{classifications}, \code{analyst\_inputs} & same & no\\
\bottomrule
\end{tabularx}

\medskip\noindent ``Required'' means that the run cannot form a view of the issuer without it. Section~\ref{sec:kinds}
gives the exact effect of each gap.

\section{The kinds}\label{sec:kinds}

\kind{Issuer identity (\code{Issuer})}
{\code{ticker}, \code{gvkey} (the key of the financials, or None), \code{cik} (the key of the filings, or None),
\code{name}.}
{None. Identity is static reference data.}
{Every step. The process calls \code{view.issuer(ticker)} first.}
{\code{DataUnavailable}. Step 1 fails, the run prints the reason, and the issuer drops out of the run.}

\kind{Quarterly financials (\code{IssuerFinancials})}
{Key: the issuer's \code{gvkey}. One record per quarter and per version. USD millions, flows for the quarter,
balances at the quarter end.
\begin{itemize}
\item Income: \code{revenue}, \code{cogs}, \code{ebitda\_reported}, \code{ebitda\_adj\_company}, \code{ebitda\_covenant},
  \code{depreciation}, \code{rent\_expense}.
\item Nonrecurring items: tuple of \code{NonRecurringItem(name, amount, recurs, cash)}.
\item Cash flows: \code{operating\_cash\_flow}, \code{capex\_total}, \code{capex\_maintenance}, \code{cash\_interest},
  \code{cash\_taxes}.
\item Working capital: \code{inventory}, \code{receivables}, \code{payables}.
\item Liquidity: \code{cash}, \code{revolver\_commitment}, \code{revolver\_drawn}, \code{borrowing\_base} (infinity when
  none), \code{letters\_of\_credit}.
\item Debt: \code{debt\_current}, \code{debt\_long\_term}, \code{maturity\_schedule} (five yearly amounts),
  \code{pension\_underfunded}, \code{guarantees\_weak}. \code{source} is provenance text.
\end{itemize}}
{\code{as\_of} is the quarter end. \code{available\_at} is the first moment that version of the quarter was public.
A restatement is a new version. The view returns every quarter in its latest version public at the decision time.}
{Framework F3 normalisation (step 3: LTM table, earnings-quality tests, margin trend, FCF). F4 liquidity and the
scenario start point (methodology 5.1). F5 coverage. G3 monitoring (step 12).}
{Step 3 runs with no quarters, so the LTM figures are empty and the memo lists open questions. The scenario start
point needs at least one quarter. With none, it raises \code{DataUnavailable}, and the process marks steps 4 to 15 of that issuer as blocked with the reason. The run continues for the other issuers.}

\subsection{Filing facts (\code{IssuerFacts})}
\noindent Fields: \code{issuer\_ticker} and \code{values}, a dict of \code{name: (value, source text)}. The source text
travels to the memo. Time stamps: \code{as\_of} and \code{available\_at} are the 10-Q acceptance time. Used by: step 1
inputs, F2 framing, methodology 4.7 clocks, F5, the scenario engine. When missing: \code{DataUnavailable} stops step 1
and the issuer drops out of the run. A missing single value is None with a source text that says why, and the memo
shows it as an open question. The names that the analyst reads:

\medskip
\noindent\begin{tabularx}{\linewidth}{@{}p{5.2cm}p{3.2cm}L@{}}
\toprule
Name & Unit & Meaning\\\midrule
\code{ticker}, \code{name}, \code{cik}, \code{gvkey} & text & identity\\
\code{q}, \code{k} & text & accession of the latest 10-Q and 10-K (\code{k} may be None)\\
\code{q\_acceptance}, \code{k\_acceptance} & time & their acceptance times\\
\code{balance\_sheet\_date} & date & the 10-Q balance-sheet date\\
\code{ltm\_ebitda}, \code{ltm\_revenue} & USD m & last-twelve-months totals\\
\code{run\_rate\_revenue} & USD m a year & latest quarter x 4\\
\code{interest\_rate} & decimal & cash interest / debt\\
\code{annual\_lease\_cost} & USD m a year & operating lease cost\\
\code{lease\_wavg\_term} & years & check only\\
\code{uncommenced\_window}, \code{uncommenced\_term} & years & signed leases not yet started\\
\code{debt\_tail\_years} & years & length of the ``thereafter'' debt bucket\\
\code{rpo\_total} & USD m & remaining performance obligation\\
\code{rpo\_windows} & (start, length, share) & timing of the RPO, years and decimal\\
\code{rpo\_tail\_years} & years & length of the RPO remainder\\
\code{purchase} & (USD m, years) & purchase commitments without a ladder\\
\code{archetype} & A1 to A7 & framework F2 archetype, from the analyst\\
\code{rating}, \code{recourse}, \code{scenario\_extra\_lease} & text or None & analyst judgement\\
\code{available\_at} & time & the latest release among the inputs\\
\bottomrule
\end{tabularx}

\kind{Filing list (\code{Filing})}
{\code{issuer\_ticker}, \code{form} (for example ``10-Q''), \code{accession}.}
{\code{as\_of} and \code{available\_at} are the acceptance time.}
{Framework F11 events (step 12 estimates the next results date from the filing cadence). The walk-forward audit.}
{An empty list. Step 12 then has no results-date event.}

\kind{Balance sheet and schedules (\code{BalanceSheet}, \code{Ladder})}
{One record per filing (the 10-Q and the 10-K of the facts). USD millions, NaN when the filing does not state it.
\begin{itemize}
\item \code{debt}, \code{debt\_source}: total financial debt and the tag or rule that gave it.
\item \code{operating\_lease\_liability}, \code{finance\_lease\_liability}, \code{cash}.
\item \code{revolver\_capacity}: unused or maximum credit-line capacity. None when the filing states none.
\item \code{ladders}: \code{operating\_lease}, \code{finance\_lease}, \code{debt}, \code{purchase}, each a \code{Ladder} or
  None. \code{Ladder.amounts} are bucket amounts (the last is ``thereafter''). \code{first\_length} is the first
  bucket's length in years. The other buckets are one year. \code{total} and \code{reconciles} check the sum.
\item \code{lease\_not\_commenced}: (total USD m, shortest term years, longest term years).
\item \code{customer\_shares}: named customer to share of revenue, decimal.
\end{itemize}}
{\code{as\_of} is the balance-sheet date. \code{available\_at} is the acceptance time.}
{Methodology 4.7 clocks (obligation ladders). F5 leverage (debt, leases, cash). The F4 funding gap and the scenario
engine (debt lines, revolver).}
{\code{DataUnavailable} for the 10-Q stops step 1 for the issuer. A missing ladder is None. Then the debt ladder comes
from the balance-sheet total spread over ten years, and the memo states that assumption.}

\kind{Filing text (\code{FilingText})}
{\code{issuer\_ticker}, \code{accession}, \code{text} (plain text of the primary document).}
{\code{as\_of} and \code{available\_at} are the acceptance time.}
{F2 ``what it sells'' (the first two sentences of Item 1). F5 covenant levels (step 5).}
{The business line is empty, and the memo says ``no filing text''. The covenant table stays open.}

\kind{First-reported cash flows (\code{ReportedCashFlow})}
{USD millions for the quarter: \code{operating\_cash\_flow}, \code{capex}, \code{interest\_paid}, \code{taxes\_paid},
and the \code{accession} that stated them.}
{\code{as\_of} is the quarter end. \code{available\_at} is the acceptance time of the first filing that reported the
quarter.}
{Step 3 only. It fills a quarter whose cash flows the financials leave blank (methodology 3.2).}
{The quarter stays blank, and the LTM uses the rule in \code{hy.state.ltm\_total}.}

\kind{Equity market value (\code{MarketValue})}
{\code{issuer\_ticker}, \code{value} (price x shares, USD millions), \code{source}.}
{\code{as\_of} is the quarter end of the price. \code{available\_at} is the report date.}
{F7.7 valuation (step 8: market TEV, debt / market TEV). Kricheff and S\&P ratios (step 5). Estimate E9.}
{The market value is NaN with the reason. Those ratios are then NaN.}

\kind{Bond terms (\code{Bond})}
{\code{cusip}, \code{issuer\_ticker}, \code{coupon} (decimal), \code{maturity} (date), \code{seniority}, \code{fixed},
\code{convertible}, \code{grade} (``H'' for high yield). The domain, not the source, keeps the fixed, non-convertible,
senior unsecured or secured bonds.}
{None. Terms are static between issue and redemption.}
{Step 1 bond curve. The F5 capital structure table. F4 refinancing (coupons of bonds due inside 27 months).}
{An empty list. Step 1 is partial with an open question. Steps 4 to 15 do not run for the issuer, because it has no
spread.}

\kind{Bond observations (frame, \code{BOND\_OBS\_COLUMNS})}
{In this order: \code{cusip}, \code{date} (trade day), \code{price} (per 100 face, volume-weighted), \code{yield}
(decimal, volume-weighted), \code{n\_trades}, \code{par\_volume} (USD), \code{available\_at}.}
{\code{date} is the trade day. \code{available\_at} is when the day's trades became public (TRACE: the next day). The
view asks for the 180 days before the decision and keeps rows public by then.}
{Step 1 latest snapshot (10 days), G-spreads and the issuer curve, the 5-year spread, the pair gap history
(methodology 7.5), the recent-issue point of F4.}
{An empty frame works like 3.9. A source error stops step 1 for the issuer.}

\kind{Treasury par curve (\code{Curve})}
{\code{tenors} (years) and \code{yields} (decimal par yields), one value per tenor. NaN marks a tenor with no quote.
Two reads: \code{curve(on)} gives the latest curve public at \code{on}. \code{curves(start, end)} gives every daily
curve in the range, which the view returns as a frame (\code{curve\_history}).}
{\code{as\_of} is the curve date. \code{available\_at} is when it was published (US Treasury: the next day).}
{G-spreads, cash and Treasury returns (step 15), and the pair gap history (step 9b).}
{The run needs the current curve and stops without it. An empty history leaves the pair gap persistence as NaN.}

\kind{HY index level (\code{IndexLevel})}
{\code{spread}: the market's high yield spread level, decimal. \code{distress}: the share of high yield bonds wider
than 1,000 bp, decimal. \code{name} says how the level was made. The one index rule is
\code{hy.state.hy\_index\_level} and the one distress rule is \code{hy.state.distress\_ratio}. The TRACE adapter
applies both to its benchmark bonds. Another source may supply an index OAS directly, for example ICE BofA H0A0.}
{\code{as\_of} is the market date. \code{available\_at} is when it was public.}
{The F6 phase vote (now, 3 months ago, and once a year for 5 years). The fair spread fallback (methodology 7.2). DTS
hedges (steps 9 and 15). Estimate E0. The macro paths.}
{A missing past snapshot is NaN in the phase table, and the vote uses the others. The current level is required.}

\kind{Macro forecasts (\code{ForecastRow}) and their status}
{\code{source}, \code{variable}, \code{horizon\_q} (quarters ahead, 0 to 8), \code{vintage}, \code{available\_at},
\code{point}, \code{q10}, \code{q50}, \code{q90}, \code{method}. Variables: \code{gdp\_growth}, \code{cpi},
\code{unemployment}, \code{policy\_rate}, \code{y2}, \code{y10}, \code{sloos\_tightening}, \code{hy\_spread},
\code{distress\_ratio}, all decimals. \code{forecast\_status} returns one row per forecast source: \code{source},
\code{status}, \code{last\_vintage}, \code{last\_available\_at}, \code{rows}, \code{rows\_after\_decision\_time}. The run
writes it to \code{\_run.json} as its audit trail.}
{\code{vintage} is the forecast date. \code{available\_at} is its release.}
{Methodology 5.4 macro paths (\code{hy.macro.build\_paths}). The F6 phase vote (SLOOS). The scenario engine.}
{A source with no rows is ``stubbed''. The path builder then uses its fallback order and logs a stated assumption.}

\subsection{Peer panels (frames)}
\noindent Every panel has an \code{available\_at} column. The view drops rows that were not public at the decision
time.

\medskip
\noindent\begin{tabularx}{\linewidth}{@{}lLp{3.6cm}@{}}
\toprule
Panel & Columns and units & Used by\\\midrule
\code{bond\_panel} & \code{date} (month end), \code{cusip}, \code{gvkey}, \code{years}, \code{g\_spread} (decimal),
\code{available\_at}. The backtest also reads \code{price}, \code{yield}, \code{coupon}, \code{maturity}, \code{mod\_dur},
\code{tsy\_yield}, \code{bid\_ask\_pts}, \code{n\_trades}, \code{par\_volume} & step 9 peers, ESP and E13 (7.2, 7.4), and the
backtests\\
\code{leverage} & \code{date}, \code{gvkey}, \code{leverage} (x), \code{ltm\_ebitda}, \code{debt} (USD m),
\code{available\_at} & peer cells and the implied rating class\\
\code{gics} & \code{gvkey}, \code{gsector} (two digits), \code{indfrom}, \code{indthru}, \code{available\_at} & the peer
sector\\
\code{ratings} & \code{gvkey}, \code{datadate}, \code{splticrm} (S\&P letter), \code{available\_at} & the actual rating
class\\
\code{annual\_ebitda} & \code{gvkey}, \code{datadate}, \code{fyear}, \code{oibdp}, \code{debt} (USD m),
\code{available\_at} & SPiDER, E15\\
\code{equity\_monthly} & \code{gvkey}, \code{datadate} (month end), \code{ret} (monthly total return, decimal),
\code{available\_at} & OneScore, E16\\
\bottomrule
\end{tabularx}

\medskip\noindent When missing: without \code{bond\_panel}, \code{leverage}, \code{gics} or \code{ratings}, step 9 has
no peer cell. The fair spread then falls back to the HY index, and the memo says so. Without \code{annual\_ebitda},
SPiDER (E15) is NaN. Without \code{equity\_monthly}, OneScore (E16) is NaN, and ESP carries its weight.

\subsection{Classifications and analyst inputs}
\noindent \code{Classification} gives the GICS sector and rating class per issuer, valid from \code{as\_of}.
\code{AnalystInput} is the analyst's YAML per issuer. These are part of the protocols, but the process does not read
them through the view today. The facts adapter applies the analyst YAML, and the peer step reads the \code{gics} and
\code{ratings} panels.

\subsection{Realised outcomes (not a new kind)}
\noindent The walk-forward scores each prediction against what happened next.
\code{hy.data.outcomes.RealisedOutcomes(source, decision\_time)} reads bond observations, curves and the HY index public
after the decision. It uses the same source methods as above. It refuses a time before the decision, and no step that
forms a prediction imports it. A test checks this.

\section{How to plug in a new source}
\begin{enumerate}
\item Write one class, for example \code{MySource}, that implements every protocol in \code{hy/data/source.py}:
  \code{IssuerSource}, \code{FilingSource}, \code{BondSource}, \code{MarketSource}, \code{EquitySource},
  \code{ReferenceSource}, \code{ForecastSource}, \code{AnalystSource} and \code{PanelSource}.
  \begin{itemize}
  \item Return the record types of \code{hy/data/records.py} with the units of section 1.
  \item Stamp every record with \code{as\_of} and \code{available\_at}. Never stamp a value earlier than its real
    release.
  \item Raise \code{DataUnavailable} for an item that you do not have. Do not return a zero for a missing amount.
  \item Keep vendor field names inside your class. \code{FakeSource} in \code{hy/data/sources.py} is a short model to
    copy.
  \end{itemize}
\item Pass it to the view: \code{DataView(MySource(...), decision\_time)}.
\item Run the process: \code{process.run(tickers, end, pairs, view=DataView(MySource(...), decision))}. A run at
  \code{end} decides at \code{end} plus one day.
\item Check the connector:
  \begin{itemize}
  \item Run \code{tests/test\_pipeline\_fake\_source.py}. It shows the shape of every kind and runs the whole analyst
    on made-up data with the network blocked.
  \item Run the suite: \code{python3 -m pytest -q}.
  \item Take a golden snapshot on the new machine (\code{scripts/golden\_check.py snapshot}), then run \code{check}
    after each change. The numbers will differ from this machine's numbers because the data differs. They must not
    move between your own runs.
  \end{itemize}
\item Set \code{HY\_DATA\_ROOT} for the derived files that stay as files: the backtest summaries
  (\code{summary\_full.json}, \code{results\_h*.parquet}, \code{peer\_signals\_*.parquet}) and the run outputs. Without
  them, step 14b uses the v0 slope of spec S14b.
\end{enumerate}

\section{What stays outside the contract}
\begin{itemize}
\item \code{scripts/backtest\_alpha.py} and \code{scripts/backtest\_clocks.py} build the warehouse panels from TRACE,
  Compustat and SEC data. They are panel builders for this warehouse. Another machine supplies the panels through
  \code{panel()} instead.
\item The pilot subsets (\code{panel\_pilot}, \code{leverage\_pilot}) are warehouse files that
  \code{hy.data.history.backtest\_panel} reads by tag.
\item \code{scripts/rv\_pair.py} keeps a WRDS query for the benchmark band of its figure. The report is display only.
\item \code{scripts/walkforward.py} reads 8-K item codes from the EDGAR cache for its event marks. They are display
  only.
\end{itemize}

\end{document}
